\subsection{The three-branch image point with a double horizontal cover}

\begin{statement}{P4R.20}{Proved}
In the remaining degree-two horizontal sector after P4.55,
suppose that the generic transverse curve along the other
image component $B$ has at least three normalization branches.
The unique image-graph cycle cannot be concentrated at one
point with three local critical-curve branches, namely the
smooth branch $A$ and two local branches of $B$.
\end{statement}
\begin{proof}
Write these two local branches as $B_+,B_-$ and the
point as $p$. Let $W_1,\ldots,W_c$ be the normalized
conductor components mapping onto $B$, with degrees
$d_i>0$ over $B^\nu$. Put
\[
                         r_B=\sum_i d_i\ge3.
\]
This is the number of generic normalization branches
of the full transverse plane curve along $B$.
The normalization conductor has components
$H,R,W_1,\ldots,W_c$ and a tree graph. Its total
point-identification contribution is $c+1$.
The horizontal section $H$ maps with degree two to
$A$, while the rational vertical curve $R$ maps
with degree one. Distinct points of $H$ are not
identified in the normalization $N$. Distinct
points of $R$ are not identified either: a
contracted vertical tree effecting such an
identification would create a cycle in the ruling
fibre. All conductor branches below are counted
on their normalization.

We first record the exact spectral bound used in
the proof. Let $E=N_{A/Y}^*$, so $\deg E=2$ by
adjunction and $c_1(Y)\cdot A=0$. On $\mathbb P(E)$
use the quotient convention, tautological class
$\xi$, and fibre class $F$. Thus
\[
          \xi^2=2,\qquad \xi F=1,\qquad F^2=0,
          \qquad K_{\mathbb P(E)}=-2\xi.
\]
The actual normal derivative maps of P4.55 identify
$H$ with the normalization of the degree-two
spectral component $\Gamma_H$, and identify the
third sheet with a section $\Gamma_R\cong A$.
Remove the common zero divisor of the cubic tensor,
of degree $b\ge0$. For an integer $t$ write
\[
       \Gamma_R=\xi-tF,\qquad
       \Gamma_H=2\xi+(t-b)F.
\]
If $\delta_H$ is the normalization colength of
$\Gamma_H$ and $n=\Gamma_H\cdot\Gamma_R\ge0$,
adjunction and the intersection product give
\[
 \begin{split}
       g(H)+\delta_H&=p_a(\Gamma_H)=1+t-b,\\
       n&=4-t-b,\\
       g(H)+\delta_H+n&=5-2b.
 \end{split}                                                   \tag{R20.1}
\]
These are the actual residual components of the
cubic, with their common zero divisor retained.
Since $g(H)\ge1$, they imply $b\le2$.

\par\medskip\noindent\textbf{The exact graph when $c\ge2$.}
Let $r_A$ be the number of normalized conductor
points over the branch $A$ at $p$. It is two if
$H\to A$ ramifies there, and three otherwise.
Let $r_+,r_-$ count the points over $B_+,B_-$.
Every $W_i$ is surjective to $B^\nu$, so
$r_+,r_-\ge c$. Let $n_p$ be the number of
normalization points of $S$ over $p$.
Their nonempty supports in $\{A,+,-\}$ pairwise
intersect: distinct local hypersurface components
intersect in a curve.

A singleton support $\{A\}$ forces every point to
carry $A$. Then $n_p\le r_A$ and the graph
contribution is at least $r_++r_-\ge2c>c+1$.
A singleton $\{+\}$ forces every point to carry
$+$ and gives contribution at least
$r_A+r_-\ge c+2$; the same applies to $\{-\}$.
Thus every support has at least two elements, and
\[
 \text{contribution at }p
       \ge\frac{r_A+r_++r_-}{2}\ge c+1.
\]
The whole global budget is $c+1$. Equality is
forced throughout:
\[
                 r_A=2,\qquad r_+=r_-=c,\qquad n_p=c+1.
                                                               \tag{R20.2}
\]
Every point over $p$ has exactly two normalized
conductor branches, with different image labels,
and all graph identifications occur at $p$.
The horizontal cover ramifies over $p$.
Each $W_i$ has precisely one point over each of
$B_+,B_-$ and has valence two; $H$ and $R$ have
valence one. The connected tree is therefore
the path
\[
                    H-W_1-W_2-\cdots-W_c-R.             \tag{R20.3}
\]
Choose the label $+$ at the $H$ endpoint. The next
point on $W_1$ has label $-$. Every internal point
has support $\{+,-\}$, so it joins that point to
the $+$ point of $W_2$. Induction gives the $+$
point on the left of every $W_i$ and its $-$
point on the right. The $R$ endpoint carries $-$.

The endpoint hypersurface factors, denoted
$f_H,f_R$, thus have supports $\{A,+\}$ and
$\{A,-\}$. The $c-1$ internal factors $f_j$
have supports $\{+,-\}$. These are all the local
irreducible hypersurface factors.

\par\medskip\noindent\textbf{The actual length-two ribbon.}
The set-theoretic intersection of $f_H$ and $f_R$
is $A$ alone. Along generic $A$, the first factor
supplies two ordinary-triple branches and the
second supplies the third, so their intersection
multiplicity is two. Their one-dimensional
complete-intersection ring is Cohen--Macaulay,
without a zero-dimensional embedded component.
A coordinate $z$ on the smooth $A$ is a
nonzerodivisor. The ring is finite free of rank
two over $\mathbb C\{z\}$, and its quotient by
$z$ has length two.

If both factors had no linear term, their
restrictions modulo $z$ would belong to
$(x,y)^2$. The Artinian quotient would contain
the independent classes $1,x,y$, contrary to
length two. At least one factor is smooth.
It cannot be $f_H$: a smooth surface is its own
normalization and supplies only one generic
preimage of $A$, whereas this factor supplies
the ramified double cover. Thus $f_R$ is smooth.
Take the actual ambient coordinate $u=f_R$.
On $u=0$, the divisor of $f_H$ is exactly twice
the smooth $A$, so in coordinates
$A=(u=v=0)$ it has restriction $v^2$ times a
unit. Weierstrass preparation in $v$, followed
by depression and using $f_H\in(u,v)^2$, gives
the exact expression
\[
                    f_H=\text{unit}\,(v^2+u^2g(u,z)).     \tag{R20.4}
\]
All changes preserve $A$. The full cubic normal
tensor is therefore, up to a nonvanishing scalar,
\[
        \left(\prod_j f_j(0,0,z)\right)
                     u\bigl(v^2+u^2g(0,z)\bigr).          \tag{R20.5}
\]
The coefficient of $uv^2$ in the last factor is
one, so the common zero order is exactly the
sum of the orders of $f_j|_A$.

\par\medskip\noindent\textbf{The pole on the smooth rational branch.}
On the smooth branch $f_R=0$, its conductor is
the restriction of the product of the other
factors. The factor $f_H$ contributes $2A$.
Each $f_j$ cuts out only $B_-$ there. This
normalization branch is already smooth, and
the pole support on its minimal resolution is
SNC. Hence $A$ and its $B_-$ curve are smooth
and transverse. Choose the latter's equation
$z=0$ on this surface, with $z$ still a
coordinate on $A$, and extend it to the ambient
threefold. Then
\[
             f_j|_{f_R=0}=\text{unit}\,z^{e_j},
                  \qquad e_j\ge1.
\]
Their sum $e=\sum e_j$ is the generic conductor
exponent of the global component $W_c$ on
this branch. By (R20.5) it is also the local
common cubic-zero order. The branch on
$f_R=0$ is smooth; the plane-curve conductor
formula therefore gives
\[
          e=\sum_{\text{other branches}}I
                         \ge r_B-1\ge2.
\]
Combining $b\ge e$ with (R20.1) forces
\[
     b=e=2,quad r_B=3,quad
     g(H)=1,quad \delta_H=n=0.                         \tag{R20.6}
\]
Thus $\Gamma_H$ is smooth and disjoint from
$\Gamma_R$, and there is no further common
zero of the cubic.

\par\medskip\noindent\textbf{The exact normalization at the ramified endpoint.}
The cover $H\to A$ ramifies at $p$ by (R20.2).
Equation (R20.4) presents its spectral cover as
$w^2+g(0,z)=0$. After translating its double
root, smoothness from (R20.6) gives
$g_z(0,0)\ne0$. Hence
\[
 \begin{split}
        N_H&=\{w^2+g(u,z)=0\},\\
        (u,w,z)&\longmapsto(u,v=uw,z)
 \end{split}                                                   \tag{R20.7}
\]
is the actual normalization near that point.
It is smooth, finite, and birational onto the
branch (R20.4); off $u=0$ its inverse is
$w=v/u$. Coordinates $(u,w)$ solve
$z=\psi(u,w)$, with $z|_H$ of order two in
$w$. The minimal normalization resolution
has no exceptional curve over this smooth
point, and $H$ is $u=0$.

The other conductor branch $W_1$ meets $H$
transversely by SNC. Thus $u$ is a coordinate
on $W_1$ and has order one on its image
$B_+$. By (R20.2), $W_1$ has only one point
over the normalized $B_+$ point, so its
global degree is $d_1=1$. The same smooth
SNC argument on $f_R=0$ gives $d_c=1$.
Since $\sum d_i=r_B=3$, the case $c=2$
would have sum two, while $c>3$ is impossible.
Therefore
\[
                 c=3,\qquad d_1=d_2=d_3=1.              \tag{R20.8}
\]
There are exactly two internal factors, and
$e_1=e_2=1$.

\par\medskip\noindent\textbf{The conductor mismatch.}
Each $f_j$ intersects $f_H$ only along $B_+$.
Its pullback to the actual smooth normalization
(R20.7) consequently has divisor supported
only on $W_1$. Its restriction to $H$ vanishes
to order two: the order on $A$ is one and
$H\to A$ has ramification index two.
Transversality of $H$ and $W_1$ makes the
divisor exactly $2W_1$. Thus the generic
intersection multiplicity of $f_H$ with
each internal factor along $B_+$ is two.

The path (R20.3) identifies the same global
component $W_3$ on the smooth $R$ branch
over $B_-$ and on an internal factor over
$B_+$. At $B_-$ its conductor exponent is
$e=2$. At $B_+$ it has intersection
multiplicity two with $f_H$ and positive
intersection multiplicity with the other
internal factor. Its conductor exponent
is at least three. These cannot be two
values of the coefficient on one global
irreducible conductor component.

The full marking is also explicit. Over
$B_+$ the factors $(f_H,f_1,f_2)$ carry
$(W_1,W_2,W_3)$; over $B_-$ the factors
$(f_1,f_2,f_R)$ carry the same ordered
global components. If $p_+,p_-\ge1$ are
the intersections of the internal factors,
the conductor vectors are
\[
         (4,2+p_+,2+p_+),\qquad
         (1+p_-,1+p_-,2).
\]
There is no missing intrinsic branch-delta
term: the exponent two of the smooth $R$
branch forces its two intersections to
equal one, which makes both other generic
$B$ branches smooth as well. The third
entries already give the contradiction.

\par\medskip\noindent\textbf{The case $c=1$.}
Now the conductor graph has $H,R,W$ and
budget two. The points of $W$ over $B_+$
and $B_-$ cannot be identified in $N$,
since that creates a loop in the conductor
tree. Points carrying the two labels are
distinct; their supports intersect, so
both carry $A$. Each contributes at least
one, exhausting the graph budget. Each
has exactly one $A$ and one $B$ branch,
and no other identification occurs.

If $H$ ramifies over $p$, these two
$A$ branches are $H$ and $R$. There are
only the factors with supports
$\{A,+\}$ and $\{A,-\}$. The length-two
ribbon argument makes $f_R$ smooth.
No other factor contains $B_-$, and this
smooth branch has no singularity there.
Thus $B_-$ cannot be a critical curve.

If $H$ is unramified, the three normalized
$A$ preimages are distinct. Two occur
at the $B$ attachments and the third
is an $A$-only normalization point.
There are three factors, each supplying
one generic $D_4$ branch. Every pair
meets only along $A$, with multiplicity
one. Its complete-intersection curve
is the reduced smooth $A$: it is pure,
generically reduced, and has no other
component. Modulo $z$, length one forces
independent normal linear terms.
All three factors are smooth near $p$
and their only pairwise intersection is
$A$. Neither $B_+$ nor $B_-$ can be a
critical branch. This final contradiction
completes the proof.
\end{proof}

The exact spectral intersections, length-two
ribbon basis, normalization equation, scalar
cubic coefficient, endpoint-degree list and
marked conductor vectors are executed in
\path{code/part4/degree_two_triple_point.py}.
Its \path{.json} and \path{.out} certificates
report PASS.
